\documentclass[10pt,a4paper]{amsart}
\usepackage{amsmath,amssymb,amsthm,graphicx}
\usepackage{verbatim}
\newenvironment{reference}{\comment}{\endcomment}
%\newenvironment{reference}{}{}
\newenvironment{slogan}{\comment}{\endcomment}
\newenvironment{history}{\comment}{\endcomment}
% For commutative diagrams we use Xy-pic
\usepackage[all]{xy}

% We use 2cell for 2-commutative diagrams.
\xyoption{2cell}
% Theorem environments.
%
\theoremstyle{plain}
\newtheorem{theorem}[subsection]{Theorem}
\newtheorem{proposition}[subsection]{Proposition}
\newtheorem{lemma}[subsection]{Lemma}

\theoremstyle{definition}
\newtheorem{definition}[subsection]{Definition}
\newtheorem{example}[subsection]{Example}
\newtheorem{exercise}[subsection]{Exercise}
\newtheorem{situation}[subsection]{Situation}

\theoremstyle{remark}
\newtheorem{remark}[subsection]{Remark}
\newtheorem{remarks}[subsection]{Remarks}

\numberwithin{equation}{subsection}

% Macros
%
\def\lim{\mathop{\mathrm{lim}}\nolimits}
\def\colim{\mathop{\mathrm{colim}}\nolimits}
\def\Spec{\mathop{\mathrm{Spec}}}
\def\Hom{\mathop{\mathrm{Hom}}\nolimits}
\def\Ext{\mathop{\mathrm{Ext}}\nolimits}
\def\SheafHom{\mathop{\mathcal{H}\!\mathit{om}}\nolimits}
\def\SheafExt{\mathop{\mathcal{E}\!\mathit{xt}}\nolimits}
\def\Sch{\mathit{Sch}}
\def\Mor{\mathop{\mathrm{Mor}}\nolimits}
\def\Ob{\mathop{\mathrm{Ob}}\nolimits}
\def\Sh{\mathop{\mathit{Sh}}\nolimits}
\def\NL{\mathop{N\!L}\nolimits}
\def\CH{\mathop{\mathrm{CH}}\nolimits}
\def\proetale{{pro\text{-}\acute{e}tale}}
\def\etale{{\acute{e}tale}}
\def\QCoh{\mathit{QCoh}}
\def\Ker{\mathop{\mathrm{Ker}}}
\def\Im{\mathop{\mathrm{Im}}}
\def\Coker{\mathop{\mathrm{Coker}}}
\def\Coim{\mathop{\mathrm{Coim}}}

% Boxtimes
%
\DeclareMathSymbol{\boxtimes}{\mathbin}{AMSa}{"02}

%
% Macros for moduli stacks/spaces
%
\def\QCohstack{\mathcal{QC}\!\mathit{oh}}
\def\Cohstack{\mathcal{C}\!\mathit{oh}}
\def\Spacesstack{\mathcal{S}\!\mathit{paces}}
\def\Quotfunctor{\mathrm{Quot}}
\def\Hilbfunctor{\mathrm{Hilb}}
\def\Curvesstack{\mathcal{C}\!\mathit{urves}}
\def\Polarizedstack{\mathcal{P}\!\mathit{olarized}}
\def\Complexesstack{\mathcal{C}\!\mathit{omplexes}}
% \Pic is the operator that assigns to X its picard group, usage \Pic(X)
% \Picardstack_{X/B} denotes the Picard stack of X over B
% \Picardfunctor_{X/B} denotes the Picard functor of X over B
\def\Pic{\mathop{\mathrm{Pic}}\nolimits}
\def\Picardstack{\mathcal{P}\!\mathit{ic}}
\def\Picardfunctor{\mathrm{Pic}}
\def\Deformationcategory{\mathcal{D}\!\mathit{ef}}
\usepackage{cleveref}
\crefname{originalreference}{original reference}{original references}
\Crefname{originalreference}{Original reference}{Original references}
\makeatletter
\newlabel{algebraic-section-stack-to-presentation}{{algebraic-section-stack-to-presentation}{source}{}{sourcepointer-1}{}}
\newlabel{algebraic-section-stack-to-presentation@cref}{{[originalreference][0][]algebraic-section-stack-to-presentation}{[1][0][]source}{}{}{}}
\makeatother
\crefname{originalreference}{original reference}{original references}
\Crefname{originalreference}{Original reference}{Original references}
\makeatletter
\newlabel{bootstrap-lemma-quotient-free-action}{{bootstrap-lemma-quotient-free-action}{source}{}{sourcepointer-2}{}}
\newlabel{bootstrap-lemma-quotient-free-action@cref}{{[originalreference][0][]bootstrap-lemma-quotient-free-action}{[1][0][]source}{}{}{}}
\makeatother
\crefname{originalreference}{original reference}{original references}
\Crefname{originalreference}{Original reference}{Original references}
\makeatletter
\newlabel{bootstrap-lemma-surjective-flat-locally-finite-presentation}{{bootstrap-lemma-surjective-flat-locally-finite-presentation}{source}{}{sourcepointer-3}{}}
\newlabel{bootstrap-lemma-surjective-flat-locally-finite-presentation@cref}{{[originalreference][0][]bootstrap-lemma-surjective-flat-locally-finite-presentation}{[1][0][]source}{}{}{}}
\makeatother
\crefname{originalreference}{original reference}{original references}
\Crefname{originalreference}{Original reference}{Original references}
\makeatletter
\newlabel{bootstrap-theorem-final-bootstrap}{{bootstrap-theorem-final-bootstrap}{source}{}{sourcepointer-4}{}}
\newlabel{bootstrap-theorem-final-bootstrap@cref}{{[originalreference][0][]bootstrap-theorem-final-bootstrap}{[1][0][]source}{}{}{}}
\makeatother
\crefname{originallemma}{original lemma}{original lemmas}
\Crefname{originallemma}{Original lemma}{Original lemmas}
\makeatletter
\newlabel{lemma-base-change-finite-presentation}{{lemma-base-change-finite-presentation}{source}{}{sourcepointer-5}{}}
\newlabel{lemma-base-change-finite-presentation@cref}{{[originallemma][0][]lemma-base-change-finite-presentation}{[1][0][]source}{}{}{}}
\makeatother
\crefname{originallemma}{original lemma}{original lemmas}
\Crefname{originallemma}{Original lemma}{Original lemmas}
\makeatletter
\newlabel{lemma-base-change-flat}{{lemma-base-change-flat}{source}{}{sourcepointer-6}{}}
\newlabel{lemma-base-change-flat@cref}{{[originallemma][0][]lemma-base-change-flat}{[1][0][]source}{}{}{}}
\makeatother
\crefname{originallemma}{original lemma}{original lemmas}
\Crefname{originallemma}{Original lemma}{Original lemmas}
\makeatletter
\newlabel{lemma-base-change-gerbe}{{lemma-base-change-gerbe}{source}{}{sourcepointer-7}{}}
\newlabel{lemma-base-change-gerbe@cref}{{[originallemma][0][]lemma-base-change-gerbe}{[1][0][]source}{}{}{}}
\makeatother
\crefname{originallemma}{original lemma}{original lemmas}
\Crefname{originallemma}{Original lemma}{Original lemmas}
\makeatletter
\newlabel{lemma-cartesian-square-inertia}{{lemma-cartesian-square-inertia}{source}{}{sourcepointer-8}{}}
\newlabel{lemma-cartesian-square-inertia@cref}{{[originallemma][0][]lemma-cartesian-square-inertia}{[1][0][]source}{}{}{}}
\makeatother
\crefname{originallemma}{original lemma}{original lemmas}
\Crefname{originallemma}{Original lemma}{Original lemmas}
\makeatletter
\newlabel{lemma-descent-finite-presentation}{{lemma-descent-finite-presentation}{source}{}{sourcepointer-9}{}}
\newlabel{lemma-descent-finite-presentation@cref}{{[originallemma][0][]lemma-descent-finite-presentation}{[1][0][]source}{}{}{}}
\makeatother
\crefname{originallemma}{original lemma}{original lemmas}
\Crefname{originallemma}{Original lemma}{Original lemmas}
\makeatletter
\newlabel{lemma-descent-flat}{{lemma-descent-flat}{source}{}{sourcepointer-10}{}}
\newlabel{lemma-descent-flat@cref}{{[originallemma][0][]lemma-descent-flat}{[1][0][]source}{}{}{}}
\makeatother
\crefname{originallemma}{original lemma}{original lemmas}
\Crefname{originallemma}{Original lemma}{Original lemmas}
\makeatletter
\newlabel{lemma-gerbe-with-section}{{lemma-gerbe-with-section}{source}{}{sourcepointer-11}{}}
\newlabel{lemma-gerbe-with-section@cref}{{[originallemma][0][]lemma-gerbe-with-section}{[1][0][]source}{}{}{}}
\makeatother
\crefname{originallemma}{original lemma}{original lemmas}
\Crefname{originallemma}{Original lemma}{Original lemmas}
\makeatletter
\newlabel{lemma-local-structure-gerbe}{{lemma-local-structure-gerbe}{source}{}{sourcepointer-12}{}}
\newlabel{lemma-local-structure-gerbe@cref}{{[originallemma][0][]lemma-local-structure-gerbe}{[1][0][]source}{}{}{}}
\makeatother
\crefname{originallemma}{original lemma}{original lemmas}
\Crefname{originallemma}{Original lemma}{Original lemmas}
\makeatletter
\newlabel{lemma-presentation-inertia}{{lemma-presentation-inertia}{source}{}{sourcepointer-13}{}}
\newlabel{lemma-presentation-inertia@cref}{{[originallemma][0][]lemma-presentation-inertia}{[1][0][]source}{}{}{}}
\makeatother
\crefname{originalreference}{original reference}{original references}
\Crefname{originalreference}{Original reference}{Original references}
\makeatletter
\newlabel{spaces-descent-lemma-locally-finite-presentation-fppf-local-source}{{spaces-descent-lemma-locally-finite-presentation-fppf-local-source}{source}{}{sourcepointer-14}{}}
\newlabel{spaces-descent-lemma-locally-finite-presentation-fppf-local-source@cref}{{[originalreference][0][]spaces-descent-lemma-locally-finite-presentation-fppf-local-source}{[1][0][]source}{}{}{}}
\makeatother
\crefname{originalreference}{original reference}{original references}
\Crefname{originalreference}{Original reference}{Original references}
\makeatletter
\newlabel{spaces-groupoids-definition-stabilizer-groupoid}{{spaces-groupoids-definition-stabilizer-groupoid}{source}{}{sourcepointer-15}{}}
\newlabel{spaces-groupoids-definition-stabilizer-groupoid@cref}{{[originalreference][0][]spaces-groupoids-definition-stabilizer-groupoid}{[1][0][]source}{}{}{}}
\makeatother
\crefname{originalreference}{original reference}{original references}
\Crefname{originalreference}{Original reference}{Original references}
\makeatletter
\newlabel{spaces-groupoids-lemma-groupoid-pre-equivalence}{{spaces-groupoids-lemma-groupoid-pre-equivalence}{source}{}{sourcepointer-16}{}}
\newlabel{spaces-groupoids-lemma-groupoid-pre-equivalence@cref}{{[originalreference][0][]spaces-groupoids-lemma-groupoid-pre-equivalence}{[1][0][]source}{}{}{}}
\makeatother
\crefname{originalreference}{original reference}{original references}
\Crefname{originalreference}{Original reference}{Original references}
\makeatletter
\newlabel{spaces-groupoids-lemma-when-gerbe}{{spaces-groupoids-lemma-when-gerbe}{source}{}{sourcepointer-17}{}}
\newlabel{spaces-groupoids-lemma-when-gerbe@cref}{{[originalreference][0][]spaces-groupoids-lemma-when-gerbe}{[1][0][]source}{}{}{}}
\makeatother
\crefname{originalreference}{original reference}{original references}
\Crefname{originalreference}{Original reference}{Original references}
\makeatletter
\newlabel{spaces-morphisms-lemma-flat-permanence}{{spaces-morphisms-lemma-flat-permanence}{source}{}{sourcepointer-18}{}}
\newlabel{spaces-morphisms-lemma-flat-permanence@cref}{{[originalreference][0][]spaces-morphisms-lemma-flat-permanence}{[1][0][]source}{}{}{}}
\makeatother
\title{A stack is a gerbe over an algebraic space exactly when its inertia is flat and locally of finite presentation}
\author{The Stacks Project Authors}
\date{stacks-project:stacks-morphisms.tex:local-snapshot-20261012; GNU FDL 1.2 or later}
\begin{document}
\maketitle
\noindent\textit{Editorial note.} All delivered mathematical text is the author's own, copied from the pinned Stacks Project chapter source sources/stacks-algebraic-stacks-morphisms/pinned-original-source.tex; no statement or proof step was rewritten, reordered or removed.\par\medskip
\noindent\textit{Editorial note.} Preamble provenance: the Stacks Project keeps one single shared preamble for every chapter; this card takes it directly from the staged sources/stacks-algebraic-stacks-morphisms/preamble.tex as three verbatim spans (lines 8-12, 14-18 and 154-217), so no derived file was written and nothing was added or redefined. The wrapper class and the base packages are supplied by the assembler.\par\medskip
\noindent\textit{Editorial note.} Bibliography provenance: neither the counted statement nor its proof contains a citation (this lane rescans the whole span before writing the card), so the card carries no bibliography and no entry was reconstructed.\par\medskip
\noindent\textit{Editorial note.} Selection provenance: this card was opened under the file-based control contract (CONTROL.md), with the mechanical pre-card gate run on its statement and proof interval before the card was written.\par\medskip
\noindent\textit{Allowed standard prerequisites (editorial).}\begin{itemize}
\item The base change permanence of the gerbe condition and the fibre product description of the inertia of a base change. Invoked inside the counted proof (source lines 6333-6338). Reduces the flatness and finite presentation of the inertia to a presentation of the stack.
\item The descent statements for flatness and for local finite presentation. Invoked twice inside the counted proof (source lines 6339-6342 and 6364-6366). Allow the properties of the inertia to be checked after a flat finitely presented base change.
\item The local structure of a gerbe as a quotient stack, with the group flat and locally of finite presentation over the base. Invoked inside the counted proof (source lines 6343-6355). Computes the inertia of a quotient stack as the acting group.
\item The smooth presentation of an algebraic stack by a groupoid, with its stabilizer group algebraic space. Used inside the counted proof (source lines 6367-6395). Realises the inertia through the stabilizer of a presentation in the converse direction.
\end{itemize}\medskip
\ifdefined\hypertarget\hypertarget{sourcepointer-1}{}\fi
\noindent\textit{Source reference.} Original reference algebraic-section-stack-to-presentation: sources/stacks-algebraic-stacks-morphisms/pinned-original-source.tex, source lines 6365--6365. Navigation pointer only: the counted text refers to the Stacks item algebraic-section-stack-to-presentation, which belongs to another chapter of the project and is not part of this card. The pointer carries the target's own label as its number because this chapter's staged snapshot does not contain the other chapter's tag table; its content is not reproduced here.\par\medskip
\ifdefined\hypertarget\hypertarget{sourcepointer-2}{}\fi
\noindent\textit{Source reference.} Original reference bootstrap-lemma-quotient-free-action: sources/stacks-algebraic-stacks-morphisms/pinned-original-source.tex, source lines 6384--6384. Navigation pointer only: the counted text refers to the Stacks item bootstrap-lemma-quotient-free-action, which belongs to another chapter of the project and is not part of this card. The pointer carries the target's own label as its number because this chapter's staged snapshot does not contain the other chapter's tag table; its content is not reproduced here.\par\medskip
\ifdefined\hypertarget\hypertarget{sourcepointer-3}{}\fi
\noindent\textit{Source reference.} Original reference bootstrap-lemma-surjective-flat-locally-finite-presentation: sources/stacks-algebraic-stacks-morphisms/pinned-original-source.tex, source lines 6424--6424. Navigation pointer only: the counted text refers to the Stacks item bootstrap-lemma-surjective-flat-locally-finite-presentation, which belongs to another chapter of the project and is not part of this card. The pointer carries the target's own label as its number because this chapter's staged snapshot does not contain the other chapter's tag table; its content is not reproduced here.\par\medskip
\ifdefined\hypertarget\hypertarget{sourcepointer-4}{}\fi
\noindent\textit{Source reference.} Original reference bootstrap-theorem-final-bootstrap: sources/stacks-algebraic-stacks-morphisms/pinned-original-source.tex, source lines 4487--4487. Navigation pointer only: the counted text refers to the Stacks item bootstrap-theorem-final-bootstrap, which belongs to another chapter of the project and is not part of this card. The pointer carries the target's own label as its number because this chapter's staged snapshot does not contain the other chapter's tag table; its content is not reproduced here.\par\medskip
\ifdefined\hypertarget\hypertarget{sourcepointer-5}{}\fi
\noindent\textit{Source reference.} Original lemma lemma-base-change-finite-presentation: sources/stacks-algebraic-stacks-morphisms/pinned-original-source.tex, source lines 5628--5628. Navigation pointer only: the counted text refers to the Stacks item lemma-base-change-finite-presentation, which is an existing labeled result of this same chapter and is not part of this card. The pointer carries the target's own label as its number and its definition line as its span; its content is not reproduced here.\par\medskip
\ifdefined\hypertarget\hypertarget{sourcepointer-6}{}\fi
\noindent\textit{Source reference.} Original lemma lemma-base-change-flat: sources/stacks-algebraic-stacks-morphisms/pinned-original-source.tex, source lines 5376--5376. Navigation pointer only: the counted text refers to the Stacks item lemma-base-change-flat, which is an existing labeled result of this same chapter and is not part of this card. The pointer carries the target's own label as its number and its definition line as its span; its content is not reproduced here.\par\medskip
\ifdefined\hypertarget\hypertarget{sourcepointer-7}{}\fi
\noindent\textit{Source reference.} Original lemma lemma-base-change-gerbe: sources/stacks-algebraic-stacks-morphisms/pinned-original-source.tex, source lines 6129--6129. Navigation pointer only: the counted text refers to the Stacks item lemma-base-change-gerbe, which is an existing labeled result of this same chapter and is not part of this card. The pointer carries the target's own label as its number and its definition line as its span; its content is not reproduced here.\par\medskip
\ifdefined\hypertarget\hypertarget{sourcepointer-8}{}\fi
\noindent\textit{Source reference.} Original lemma lemma-cartesian-square-inertia: sources/stacks-algebraic-stacks-morphisms/pinned-original-source.tex, source lines 1436--1436. Navigation pointer only: the counted text refers to the Stacks item lemma-cartesian-square-inertia, which is an existing labeled result of this same chapter and is not part of this card. The pointer carries the target's own label as its number and its definition line as its span; its content is not reproduced here.\par\medskip
\ifdefined\hypertarget\hypertarget{sourcepointer-9}{}\fi
\noindent\textit{Source reference.} Original lemma lemma-descent-finite-presentation: sources/stacks-algebraic-stacks-morphisms/pinned-original-source.tex, source lines 5833--5833. Navigation pointer only: the counted text refers to the Stacks item lemma-descent-finite-presentation, which is an existing labeled result of this same chapter and is not part of this card. The pointer carries the target's own label as its number and its definition line as its span; its content is not reproduced here.\par\medskip
\ifdefined\hypertarget\hypertarget{sourcepointer-10}{}\fi
\noindent\textit{Source reference.} Original lemma lemma-descent-flat: sources/stacks-algebraic-stacks-morphisms/pinned-original-source.tex, source lines 5389--5389. Navigation pointer only: the counted text refers to the Stacks item lemma-descent-flat, which is an existing labeled result of this same chapter and is not part of this card. The pointer carries the target's own label as its number and its definition line as its span; its content is not reproduced here.\par\medskip
\ifdefined\hypertarget\hypertarget{sourcepointer-11}{}\fi
\noindent\textit{Source reference.} Original lemma lemma-gerbe-with-section: sources/stacks-algebraic-stacks-morphisms/pinned-original-source.tex, source lines 6183--6183. Navigation pointer only: the counted text refers to the Stacks item lemma-gerbe-with-section, which is an existing labeled result of this same chapter and is not part of this card. The pointer carries the target's own label as its number and its definition line as its span; its content is not reproduced here.\par\medskip
\ifdefined\hypertarget\hypertarget{sourcepointer-12}{}\fi
\noindent\textit{Source reference.} Original lemma lemma-local-structure-gerbe: sources/stacks-algebraic-stacks-morphisms/pinned-original-source.tex, source lines 6248--6248. Navigation pointer only: the counted text refers to the Stacks item lemma-local-structure-gerbe, which is an existing labeled result of this same chapter and is not part of this card. The pointer carries the target's own label as its number and its definition line as its span; its content is not reproduced here.\par\medskip
\ifdefined\hypertarget\hypertarget{sourcepointer-13}{}\fi
\noindent\textit{Source reference.} Original lemma lemma-presentation-inertia: sources/stacks-algebraic-stacks-morphisms/pinned-original-source.tex, source lines 1513--1513. Navigation pointer only: the counted text refers to the Stacks item lemma-presentation-inertia, which is an existing labeled result of this same chapter and is not part of this card. The pointer carries the target's own label as its number and its definition line as its span; its content is not reproduced here.\par\medskip
\ifdefined\hypertarget\hypertarget{sourcepointer-14}{}\fi
\noindent\textit{Source reference.} Original reference spaces-descent-lemma-locally-finite-presentation-fppf-local-source: sources/stacks-algebraic-stacks-morphisms/pinned-original-source.tex, source lines 3896--3896. Navigation pointer only: the counted text refers to the Stacks item spaces-descent-lemma-locally-finite-presentation-fppf-local-source, which belongs to another chapter of the project and is not part of this card. The pointer carries the target's own label as its number because this chapter's staged snapshot does not contain the other chapter's tag table; its content is not reproduced here.\par\medskip
\ifdefined\hypertarget\hypertarget{sourcepointer-15}{}\fi
\noindent\textit{Source reference.} Original reference spaces-groupoids-definition-stabilizer-groupoid: sources/stacks-algebraic-stacks-morphisms/pinned-original-source.tex, source lines 4323--4323. Navigation pointer only: the counted text refers to the Stacks item spaces-groupoids-definition-stabilizer-groupoid, which belongs to another chapter of the project and is not part of this card. The pointer carries the target's own label as its number because this chapter's staged snapshot does not contain the other chapter's tag table; its content is not reproduced here.\par\medskip
\ifdefined\hypertarget\hypertarget{sourcepointer-16}{}\fi
\noindent\textit{Source reference.} Original reference spaces-groupoids-lemma-groupoid-pre-equivalence: sources/stacks-algebraic-stacks-morphisms/pinned-original-source.tex, source lines 6409--6409. Navigation pointer only: the counted text refers to the Stacks item spaces-groupoids-lemma-groupoid-pre-equivalence, which belongs to another chapter of the project and is not part of this card. The pointer carries the target's own label as its number because this chapter's staged snapshot does not contain the other chapter's tag table; its content is not reproduced here.\par\medskip
\ifdefined\hypertarget\hypertarget{sourcepointer-17}{}\fi
\noindent\textit{Source reference.} Original reference spaces-groupoids-lemma-when-gerbe: sources/stacks-algebraic-stacks-morphisms/pinned-original-source.tex, source lines 6420--6420. Navigation pointer only: the counted text refers to the Stacks item spaces-groupoids-lemma-when-gerbe, which belongs to another chapter of the project and is not part of this card. The pointer carries the target's own label as its number because this chapter's staged snapshot does not contain the other chapter's tag table; its content is not reproduced here.\par\medskip
\ifdefined\hypertarget\hypertarget{sourcepointer-18}{}\fi
\noindent\textit{Source reference.} Original reference spaces-morphisms-lemma-flat-permanence: sources/stacks-algebraic-stacks-morphisms/pinned-original-source.tex, source lines 5451--5451. Navigation pointer only: the counted text refers to the Stacks item spaces-morphisms-lemma-flat-permanence, which belongs to another chapter of the project and is not part of this card. The pointer carries the target's own label as its number because this chapter's staged snapshot does not contain the other chapter's tag table; its content is not reproduced here.\par\medskip
\begin{proposition}
\label{proposition-when-gerbe}
Let $\mathcal{X}$ be an algebraic stack. The following are equivalent
\begin{enumerate}
\item $\mathcal{X}$ is a gerbe, and
\item $\mathcal{I}_\mathcal{X} \to \mathcal{X}$ is flat and locally of
finite presentation.
\end{enumerate}
\end{proposition}

\begin{proof}
Assume (1). Choose a morphism $\mathcal{X} \to X$ into an algebraic space $X$
which turns $\mathcal{X}$ into a gerbe over $X$. Let $X' \to X$ be a
surjective, flat, locally finitely presented morphism and
set $\mathcal{X}' = X' \times_X \mathcal{X}$. Note that $\mathcal{X}'$
is a gerbe over $X'$ by
Lemma \ref{lemma-base-change-gerbe}.
Then both squares in
$$
\xymatrix{
\mathcal{I}_{\mathcal{X}'} \ar[r] \ar[d] &
\mathcal{X}' \ar[r] \ar[d] & X' \ar[d] \\
\mathcal{I}_\mathcal{X} \ar[r] &
\mathcal{X} \ar[r] & X
}
$$
are fibre product squares, see
Lemma \ref{lemma-cartesian-square-inertia}.
Hence to prove $\mathcal{I}_\mathcal{X} \to \mathcal{X}$ is flat and
locally of finite presentation it suffices to do so after such a base
change by
Lemmas \ref{lemma-descent-flat} and
\ref{lemma-descent-finite-presentation}.
Thus we can apply
Lemma \ref{lemma-local-structure-gerbe}
to assume that $\mathcal{X} = [U/G]$.
By
Lemma \ref{lemma-gerbe-with-section}
we see $G$ is flat and locally of finite presentation over $U$ and
that $x : U \to [U/G]$ is surjective, flat, and locally of finite
presentation. Moreover, the pullback of $\mathcal{I}_\mathcal{X}$
by $x$ is $G$ and we conclude that (2) holds by descent again, i.e., by
Lemmas \ref{lemma-descent-flat} and
\ref{lemma-descent-finite-presentation}.

\medskip\noindent
Conversely, assume (2). Choose a smooth presentation $\mathcal{X} = [U/R]$, see
Algebraic Stacks, Section \ref{algebraic-section-stack-to-presentation}.
Denote $G \to U$ the stabilizer group algebraic space of the groupoid
$(U, R, s, t, c, e, i)$, see
Groupoids in Spaces, Definition
\ref{spaces-groupoids-definition-stabilizer-groupoid}.
By
Lemma \ref{lemma-presentation-inertia}
we see that $G \to U$ is flat and locally of finite presentation as
a base change of $\mathcal{I}_\mathcal{X} \to \mathcal{X}$, see
Lemmas \ref{lemma-base-change-flat} and
\ref{lemma-base-change-finite-presentation}.
Consider the following action
$$
a : G \times_{U, t} R \to R, \quad (g, r) \mapsto c(g, r)
$$
of $G$ on $R$. This action is free on $T$-valued points for any
scheme $T$ as $R$ is a groupoid. Hence $R' = R/G$ is an algebraic
space and the quotient morphism $\pi : R \to R'$ is surjective,
flat, and locally of finite presentation by
Bootstrap, Lemma \ref{bootstrap-lemma-quotient-free-action}.
The projections $s, t : R \to U$ are $G$-invariant, hence
we obtain morphisms $s' , t' : R' \to U$ such that $s = s' \circ \pi$
and $t = t' \circ \pi$.
Since $s, t : R \to U$ are flat and locally of finite presentation
we conclude that $s', t'$ are flat and locally of finite presentation, see
Morphisms of Spaces, Lemmas
\ref{spaces-morphisms-lemma-flat-permanence} and
Descent on Spaces, Lemma
\ref{spaces-descent-lemma-locally-finite-presentation-fppf-local-source}.
Consider the morphism
$$
j' = (t', s') : R' \longrightarrow U \times U.
$$
We claim this is a monomorphism. Namely, suppose that $T$ is a scheme
and that $a, b : T \to R'$ are morphisms which have the same image
in $U \times U$. By definition of the quotient $R' = R/G$ there
exists an fppf covering $\{h_j : T_j \to T\}$ such
that $a \circ h_j = \pi \circ a_j$ and $b \circ h_j = \pi \circ b_j$
for some morphisms $a_j, b_j : T_j \to R$. Since $a_j, b_j$ have the same
image in $U \times U$ we see that $g_j = c(a_j, i(b_j))$ is a $T_j$-valued
point of $G$ such that $c(g_j, b_j) = a_j$. In other words, $a_j$ and
$b_j$ have the same image in $R'$ and the claim is proved.
Since $j : R \to U \times U$ is a pre-equivalence relation (see
Groupoids in Spaces, Lemma
\ref{spaces-groupoids-lemma-groupoid-pre-equivalence})
and $R \to R'$ is surjective (as a map of sheaves) we see that
$j' : R' \to U \times U$ is an equivalence relation.
Hence
Bootstrap, Theorem \ref{bootstrap-theorem-final-bootstrap}
shows that $X = U/R'$ is an algebraic space.
Finally, we claim that the morphism
$$
\mathcal{X} = [U/R] \longrightarrow X = U/R'
$$
turns $\mathcal{X}$ into a gerbe over $X$. This follows from
Groupoids in Spaces, Lemma \ref{spaces-groupoids-lemma-when-gerbe}
as $R \to R'$ is surjective, flat, and locally of finite presentation
(if needed use
Bootstrap, Lemma
\ref{bootstrap-lemma-surjective-flat-locally-finite-presentation}
to see this implies the required hypothesis).
\end{proof}

\end{document}
